% Options for packages loaded elsewhere
\PassOptionsToPackage{unicode}{hyperref}
\PassOptionsToPackage{hyphens}{url}
\PassOptionsToPackage{dvipsnames,svgnames,x11names}{xcolor}
\documentclass[
  10pt,
  fontset=fandol]{ctexart}
\usepackage{xcolor}
\usepackage[margin=16mm]{geometry}
\usepackage{amsmath,amssymb}
\setcounter{secnumdepth}{-\maxdimen} % remove section numbering
\usepackage{iftex}
\ifPDFTeX
  \usepackage[T1]{fontenc}
  \usepackage[utf8]{inputenc}
  \usepackage{textcomp} % provide euro and other symbols
\else % if luatex or xetex
  \usepackage{unicode-math} % this also loads fontspec
  \defaultfontfeatures{Scale=MatchLowercase}
  \defaultfontfeatures[\rmfamily]{Ligatures=TeX,Scale=1}
\fi
\usepackage{lmodern}
\ifPDFTeX\else
  % xetex/luatex font selection
\fi
% Use upquote if available, for straight quotes in verbatim environments
\IfFileExists{upquote.sty}{\usepackage{upquote}}{}
\IfFileExists{microtype.sty}{% use microtype if available
  \usepackage[]{microtype}
  \UseMicrotypeSet[protrusion]{basicmath} % disable protrusion for tt fonts
}{}
\makeatletter
\@ifundefined{KOMAClassName}{% if non-KOMA class
  \IfFileExists{parskip.sty}{%
    \usepackage{parskip}
  }{% else
    \setlength{\parindent}{0pt}
    \setlength{\parskip}{6pt plus 2pt minus 1pt}}
}{% if KOMA class
  \KOMAoptions{parskip=half}}
\makeatother
\setlength{\emergencystretch}{3em} % prevent overfull lines
\providecommand{\tightlist}{%
  \setlength{\itemsep}{0pt}\setlength{\parskip}{0pt}}
\usepackage{amsmath,amssymb,mathtools,needspace}
\xeCJKDeclareCharClass{CJK}{"2032,"0400->"04FF,"2160->"216B,"2460->"2473,"25A0,"25C7,"25CB,"25CF}
\IfFontExistsTF{Arial Unicode MS}{\xeCJKsetup{AutoFallBack=true}\setCJKfallbackfamilyfont{\CJKrmdefault}{Arial Unicode MS}}{}
\setcounter{secnumdepth}{0}
\setlength{\emergencystretch}{2em}
\allowdisplaybreaks
\usepackage{bookmark}
\IfFileExists{xurl.sty}{\usepackage{xurl}}{} % add URL line breaks if available
\urlstyle{same}
\hypersetup{
  pdftitle={最佳平方三角逼近与三角插值},
  colorlinks=true,
  linkcolor={Maroon},
  filecolor={Maroon},
  citecolor={Blue},
  urlcolor={Blue},
  pdfcreator={LaTeX via pandoc}}

\title{最佳平方三角逼近与三角插值}
\author{}
\date{}

\begin{document}
\maketitle

\subsubsection{3.7.1 最佳平方三角逼近与三角插值}\label{ux6700ux4f73ux5e73ux65b9ux4e09ux89d2ux903cux8fd1ux4e0eux4e09ux89d2ux63d2ux503c}

设 \(f(x)\) 是以 \(2\pi\) 为周期的平方可积函数，用三角多项式

\[
S_n(x)=\frac12a_0+a_1\cos x+b_1\sin x+\cdots+a_n\cos nx+b_n\sin nx\tag{3.7.1}
\]

作最佳平方逼近函数。由于三角函数族 \(1,\allowbreak \cos x,\allowbreak \sin x,\allowbreak \cdots,\allowbreak \cos kx,\allowbreak \sin kx,\allowbreak \cdots\) 在 \([0,2\pi]\) 上是正交函数族，于是 \(f(x)\) 在 \([0,2\pi]\) 上的最小平方三角逼近多项式 \(S_n(x)\) 的系数是

\begin{quote}
校注（agent 补充）：3.6.3 原书的 \(S_n\) 求和从 \(k=1\) 开始，但后续使用 \(F(a_0,a_1,\ldots,a_n)\) 和 \(a_k\;(k=0,\ldots,n)\)；段中系数列表印为 \(a_0,a_2,\ldots,a_n\)；内积公式右端第二因子的原印下标为 \(i\)，而左端为 \(j\)，相容写法应使用 \(\varphi_j\)。这些原书记号不一致均已保留，见 \href{../assets/ch03b/p084-multivariate-formulas.png}{原文裁切}。
\end{quote}

\href{../../source-images/pdf-085.jpeg}{查看原页}

\[
\begin{cases}
a_k=\dfrac1\pi\displaystyle\int_0^{2\pi}f(x)\cos kx\,\mathrm dx\quad(k=0,1,\cdots,n),\\[4pt]
b_k=\dfrac1\pi\displaystyle\int_0^{2\pi}f(x)\sin kx\,\mathrm dx\quad(k=1,2,\cdots,n),
\end{cases}\tag{3.7.2}
\]

\(a_k,b_k\) 称为 Fourier 系数，函数 \(f(x)\) 按 Fourier 系数展开得到的级数

\[
\frac12a_0+\sum_{k=1}^{\infty}(a_k\cos kx+b_k\sin kx)\tag{3.7.3}
\]

就称为 Fourier 级数，只要 \(f'(x)\) 在 \([0,2\pi]\) 上分段连续，则级数 (3.7.3) 一致收敛到 \(f(x)\)。

当 \(f(x)\) 只在给定的离散点集 \(\left\{x_j=\dfrac{2\pi}{N}j,\ j=0,1,\cdots,N-1\right\}\) 上已知时，则可类似得到在离散点集上的正交性与相应的离散 Fourier 系数。为方便起见，下面只给出奇数个点的情形。令

\[
x_j=\frac{2\pi j}{2m+1}\quad(j=0,1,\cdots,2m),
\]

可以证明，对于任何 \(0\leqslant k,l\leqslant m\)，成立

\[
\left\{\begin{aligned}
\sum_{j=0}^{2m}\sin lx_j\sin kx_j&=\begin{cases}
0,&l\ne k,\ l=k=0,\\
\dfrac{2m+1}{2},&l=k\ne0,
\end{cases}\\[4pt]
\sum_{j=0}^{2m}\cos lx_j\cos kx_j&=\begin{cases}
0,&l\ne k,\\
\dfrac{2m+1}{2},&l=k\ne0,\\
2m+1,&l=k=0,
\end{cases}\\[4pt]
\sum_{j=0}^{2m}\cos lx_j\sin kx_j&=0,\quad0\leqslant k,j\leqslant m.
\end{aligned}\right.
\]

这就表明，函数族 \(\{1,\allowbreak \cos x,\allowbreak \sin x,\allowbreak \cdots,\allowbreak \cos mx,\allowbreak \sin mx\}\) 在点集 \(\left\{x_j=\dfrac{2\pi j}{2m+1}\right\}\) 上正交。若令 \(f_j=f(x_j)\)（\(j=0,1,\cdots,2m\)），则 \(f(x)\) 的最小二乘三角逼近为

\[
S_n(x)=\frac12a_0+\sum_{k=1}^n(a_k\cos kx+b_k\sin kx),\quad n<m,
\]

其中

\[
\begin{cases}
a_k=\dfrac2{2m+1}\displaystyle\sum_{j=0}^{2m}f_j\cos\dfrac{2\pi jk}{2m+1}\quad(k=0,1,\cdots,n),\\[4pt]
b_k=\dfrac2{2m+1}\displaystyle\sum_{j=0}^{2m}f_j\sin\dfrac{2\pi jk}{2m+1}\quad(k=1,\cdots,n).
\end{cases}\tag{3.7.4}
\]

当 \(n=m\) 时，可证明 \(S_m(x_j)=f_j\)（\(j=0,1,\cdots,2m\)），于是

\[
S_m(x)=\frac12a_0+\sum_{k=1}^m(a_k\cos kx+b_k\sin kx)
\]

\begin{quote}
校注（agent 补充）：离散正交性公式的最后一行，原书条件为 \(0\leqslant k,j\leqslant m\)，已保留。这里 \(j\) 是求和指标，函数频率为 \(k,l\)；与该组公式前的条件一致时应写 \(0\leqslant k,l\leqslant m\)。
\end{quote}

\href{../../source-images/pdf-086.jpeg}{查看原页}

就是三角插值多项式，系数仍由式 (3.7.4) 表示。

更一般情形，假定 \(f(x)\) 是以 \(2\pi\) 为周期的复函数，给定在 \(N\) 个等分点 \(x_j=\dfrac{2\pi}{N}j\)（\(j=0,1,\cdots,N-1\)）上的值 \(f_j=f\left(\dfrac{2\pi}{N}j\right)\)，由于

\[
\mathrm e^{\mathrm ijx}=\cos(jx)+\mathrm i\sin(jx)\quad(j=0,1,\cdots,N-1;\ \mathrm i=\sqrt{-1}),
\]

函数族 \(\{1,\mathrm e^{\mathrm ix},\cdots,\mathrm e^{\mathrm i(N-1)x}\}\) 在区间 \([0,2\pi]\) 上是正交的。将函数 \(\mathrm e^{\mathrm ijx}\) 在等距点集 \(x_k=\dfrac{2\pi}{N}k\)（\(k=0,1,\cdots,N-1\)）上的值 \(\mathrm e^{\mathrm ijx_k}\) 组成的向量记作

\[
\boldsymbol\varphi_j=\left(1,\mathrm e^{\mathrm ij\frac{2\pi}{N}},\cdots,\mathrm e^{\mathrm ij\frac{2\pi}{N}(N-1)}\right)^{\mathrm T}.
\]

当 \(j=0,1,\cdots,N-1\) 时，\(N\) 个复向量 \(\boldsymbol\varphi_0,\boldsymbol\varphi_1,\cdots,\boldsymbol\varphi_{N-1}\) 具有下面所定义的正交性：

\[
(\boldsymbol\varphi_l,\boldsymbol\varphi_s)=\sum_{k=0}^{N-1}\mathrm e^{\mathrm il\frac{2\pi}{N}k}\mathrm e^{-\mathrm is\frac{2\pi}{N}k}=\sum_{k=0}^{N-1}\mathrm e^{\mathrm i(l-s)\frac{2\pi}{N}k}=\begin{cases}0,&l\ne s,\\N,&l=s.\end{cases}\tag{3.7.5}
\]

事实上，令 \(r=\mathrm e^{\mathrm i(l-s)\frac{2\pi}{N}}\)，若 \(0\leqslant l,s\leqslant N-1\)，则有

\[
0\leqslant l\leqslant N-1,\quad-(N-1)\leqslant-s\leqslant0,
\]

于是

\[
-(N-1)\leqslant l-s\leqslant N-1,\quad-1<-\frac{N-1}{N}\leqslant\frac{l-s}{N}\leqslant\frac{N-1}{N}<1.
\]

若 \(l-s\ne0\)，则 \(r\ne1\)，从而 \(r^N=\mathrm e^{\mathrm i(l-s)2\pi}=1\)，于是

\[
(\boldsymbol\varphi_l,\boldsymbol\varphi_s)=\sum_{k=0}^{N-1}r^k=\frac{1-r^N}{1-r}=0.
\]

若 \(l=s\)，则 \(r=1\)，于是 \((\boldsymbol\varphi_s,\boldsymbol\varphi_s)=\sum_{k=0}^{N-1}r^k=N\)。这就证明了式 (3.7.5) 成立，即 \(\boldsymbol\varphi_0,\boldsymbol\varphi_1,\cdots,\boldsymbol\varphi_{N-1}\) 是正交的。

因此，\(f(x)\) 在 \(N\) 个点 \(\left\{x_j=\dfrac{2\pi}{N}j,\ j=0,1,\cdots,N-1\right\}\) 上的最小二乘 Fourier 逼近为

\[
S(x)=\sum_{k=0}^{n-1}C_k\mathrm e^{\mathrm ikx},\quad n\leqslant N,\tag{3.7.6}
\]

其中

\[
C_k=\frac1N\sum_{j=0}^{N-1}f_j\mathrm e^{-\mathrm ikj\frac{2\pi}{N}}\quad(k=0,1,\cdots,N-1).\tag{3.7.7}
\]

在式 (3.7.6) 中，若 \(n=N\)，则 \(S(x)\) 为 \(f(x)\) 在点 \(x_j\)（\(j=0,1,\cdots,N-1\)）上的插值函数，即 \(S(x_j)=f(x_j)\)，于是由式 (3.7.6) 得

\[
f_j=\sum_{k=0}^{N-1}C_k\mathrm e^{\mathrm ik\frac{2\pi}{N}j}\quad(j=0,1,\cdots,N-1).\tag{3.7.8}
\]

式 (3.7.7) 表示了由 \(\{f_j\}\) 求 \(\{C_k\}\) 的过程，称为 \(f(x)\) 的离散 Fourier 变换，简称

\href{../../source-images/pdf-087.jpeg}{查看原页}

DFT；而式 (3.7.8) 是由 \(\{C_k\}\) 求 \(\{f_j\}\) 的过程，称为反变换。它们是使用计算机进行 Fourier 分析的主要方法，在数字信号处理、全息技术、光谱和声谱分析等很多领域都有广泛应用。

\begin{center}\rule{0.5\linewidth}{0.5pt}\end{center}

\subsection{相关原页校注（agent 补充）}\label{ux76f8ux5173ux539fux9875ux6821ux6ce8agent-ux8865ux5145}

以下校注来自本节覆盖的原页，可能也涉及同页相邻小节；原文照录与校正说明分开保留。

\subsubsection{原书 PDF 页 87 的校注}\label{ux539fux4e66-pdf-ux9875-87-ux7684ux6821ux6ce8}

\href{../pages/pdf-087.md}{完整原页转录} · \href{../../source-images/pdf-087.jpeg}{原页图像}

校注记录：ch03b-E015, ch03b-E016。

\begin{quote}
校注（agent 补充）：原书``特别当 \(N=2^p\) 时，只有 \(N/2\) 个不同的值''已保留。对原文列出的全部 \(\exp(\mathrm i2\pi kj/N)\) 而言，取 \(j=1\) 已有 \(N\) 个不同值；若利用 \(w^{r+N/2}=-w^r\)，才可用 \(N/2\) 个值及符号表示它们。因此该处``不同的值''疑漏了``除符号外''一类限定。原书``可用 \(w\) 的 \(N\) 同余数 \(r\) 代替 \(m\)''也按原文保留，根据前段定义，此处应是指数 \(m\) 的余数。
\end{quote}

\subsection{本节覆盖原页的校注记录}\label{ux672cux8282ux8986ux76d6ux539fux9875ux7684ux6821ux6ce8ux8bb0ux5f55}

下列记录按原页关联，含同页相邻小节的校注；计算时同时读取上文校注和完整原页。

\begin{itemize}
\tightlist
\item
  \texttt{ch03b-E012}：原书 PDF 页 84
\item
  \texttt{ch03b-E013}：原书 PDF 页 84
\item
  \texttt{ch03b-E014}：原书 PDF 页 85
\item
  \texttt{ch03b-E015}：原书 PDF 页 87
\item
  \texttt{ch03b-E016}：原书 PDF 页 87
\end{itemize}

\end{document}
